\textcolor{black}{\section{Introducción}}
\noindent
Se puede determinar la relación carga-masa del electrón $e/m_e$ haciendo uso de la emisión termoiónica, poder determinar esta relación experimentalmente es importante pues nos acercamos a determinar las propias constantes fundamentales.

En esta práctica se emplea un diodo de vacío de geometría cilíndrica, en el cual los electrones son emitidos desde un cátodo calentado mediante el fenómeno de emisión termoiónica. Cuando el cátodo alcanza una temperatura suficientemente elevada, una fracción de los electrones adquiere la energía necesaria para superar la función de trabajo del material y abandonar su superficie, formando una nube electrónica en el espacio comprendido entre el cátodo y el ánodo. Esta nube de electrones produce un campo eléctrico propio que modifica el movimiento de las partículas emitidas, dando origen al fenómeno llamado saturación del voltaje de la placa (plate voltage saturation).
 

Cuando el suministro de electrones por parte del cátodo es suficiente, la corriente que circula entre los electrodos deja de depender de la capacidad emisora del material y pasa a estar gobernada por la distribución del potencial eléctrico generada por la propia carga espacial. En estas condiciones, Child y Langmuir demostraron que la corriente presenta una dependencia proporcional a la potencia $3/2$ del voltaje aplicado entre el ánodo y el cátodo. Para un diodo con electrodos cilíndricos concéntricos, la ley de Child-Langmuir establece que:

\begin{equation}
	I=
	B
	\sqrt{\frac{e}{m_e}}
	\,V^{3/2},
\end{equation}
con:
 \begin{equation}
 	B = \frac{8\pi\varepsilon_0L \sqrt{2}}{9r_a\beta^2}
 \end{equation}
 
donde $I$ es la corriente del ánodo, $V$ es la diferencia de potencial entre el ánodo y el cátodo, $L$ es la longitud efectiva del cátodo, $r_a$ es el radio del ánodo, $\varepsilon_0$ es la permitividad del vacío y $\beta$ es un factor de corrección que depende de la relación entre los radios del ánodo y del cátodo. Esta expresión muestra que, conociendo los parámetros geométricos del diodo y midiendo experimentalmente la relación entre la corriente y el voltaje, es posible determinar la relación carga-masa del electrón y también verificar la dependencia $I\propto V^{3/2}$ predicha por la Ley de Child-Langmuir.

Siguiendo la Ley de Child-Langmuir en escala loglog:
\begin{equation}
	ln(I) = ln\left( B \sqrt{\frac{e}{m_e}} \right) + \frac{3}{2}ln(V)
	\label{eq:loglog}
\end{equation}

\noindent 
